\begin{enumerate}
    \item Trouver la vitesse angulaire de la grande aiguille de l’horloge
          $\omega_{G}$.
    \item Trouver la vitesse angulaire de la petite aiguille de l’horloge
          $\omega_{P}$.
    \item On considère que 12:00 est l'origine du temps. À quel instant les deux
          aiguilles seront confondus pour la première fois.
    \item Sachant que la longueur de la grande aiguille est
          $L_{1}=\qty{20}{\cm}$ et la longueur de la petite aiguille est
          $L_{2}=\qty{15}{\cm}$.
          
          Représenter les deux aiguilles et les vecteurs vitesses à leurs
          extrémités a 5h exacte.
          
          Utilisons les deux échelles~: $\qty{1}{\cm} \to \qty{8}{\cm}$
          et $\qty{1}{\cm} \to \qty{5e-5}{\v}$.
\end{enumerate}

